\documentclass[11pt]{article}

\usepackage[margin=1in]{geometry}
\usepackage{amsmath,amssymb}
\usepackage{enumitem}
\usepackage{hyperref}
\usepackage{xcolor}
\usepackage{booktabs}
\usepackage{listings}

\hypersetup{colorlinks=true, linkcolor=blue, urlcolor=blue, citecolor=blue}

\lstset{
  basicstyle=\ttfamily\small,
  breaklines=true,
  columns=fullflexible,
}

\title{TOAE209/TOAH209 -- Design and Analysis of Algorithms\\[4pt]
\large Week 13: Practical Exercises}
\author{Foreign Trade University}
\date{}

\begin{document}
\maketitle

\section*{Instructions}

For each problem below there is a starter file
\texttt{exercises/starter\_code\_problemNN.py} containing function signatures
with \texttt{raise NotImplementedError} bodies, and a small \texttt{\_\_main\_\_}
block with tests. Implement the missing functions so the tests pass. A reference
solution is provided in \texttt{solutions/solution\_problemNN.py} -- try the
problem yourself before looking!

All problems use only the Python 3.10+ standard library. Shared helper functions
and data types live in \texttt{starter\_code.py} at the week root:
\texttt{generate\_graph}, \texttt{generate\_tree}, \texttt{edges\_to\_adj},
\texttt{is\_vertex\_cover}, \texttt{is\_independent\_set},
\texttt{brute\_force\_vertex\_cover}, \texttt{brute\_force\_independent\_set},
the QBF node classes (\texttt{Const}, \texttt{Var}, \texttt{Not}, \texttt{And},
\texttt{Or}, \texttt{Quant}), \texttt{cnf\_is\_satisfied}, and
\texttt{generate\_2sat}.

The 12 problems are grouped into three themes:
\begin{itemize}
    \item \textbf{PSPACE \& Games} (Problems 1--3): recursive evaluation of
    quantified boolean formulas (the canonical PSPACE-complete problem), minimax
    game-tree evaluation, and alpha-beta pruning.
    \item \textbf{Fixed-Parameter Tractability \& Special Structure} (Problems
    4--7, 10, 12): bounded-search-tree Vertex Cover, verification against brute
    force, kernelization, maximum independent set on trees, 2-SAT via SCCs, and
    FPT set cover.
    \item \textbf{Branch and Bound \& State-Space Search} (Problems 8, 9, 11):
    branch-and-bound knapsack and TSP, and iterative-deepening DFS.
\end{itemize}

\newpage
\section{PSPACE and Games}

\subsection*{Problem 1 -- Evaluate a Fully-Quantified Boolean Formula (QBF)}

Deciding a quantified boolean formula (QSAT/TQBF) is the canonical
PSPACE-complete problem. Implement \texttt{eval\_qbf(formula, env)} that
recursively evaluates a quantified boolean formula and returns \texttt{True} or
\texttt{False}. The formula is built from the node classes in
\texttt{starter\_code.py}: \texttt{Const}, \texttt{Var}, \texttt{Not},
\texttt{And}, \texttt{Or}, and \texttt{Quant}.

A \texttt{Quant('E', var, body)} means ``$\exists$ var. body'' (true if body is
true for var $=$ False \emph{or} var $=$ True); \texttt{Quant('A', var, body)}
means ``$\forall$ var. body'' (true only if body is true for \emph{both}
values). Evaluate by trying both values, recursing into the body, and combining
with \texttt{any(...)} for $\exists$ and \texttt{all(...)} for $\forall$.

\textbf{Example.} \texttt{eval\_qbf(Quant("A", "x", Or(Var("x"),
Not(Var("x")))))} is \texttt{True} (the tautology $\forall x.\,(x \lor \lnot
x)$), while \texttt{eval\_qbf(Quant("A", "x", Var("x")))} is \texttt{False}.

\subsection*{Problem 2 -- Minimax Game-Tree Evaluation}

Implement \texttt{minimax(node, maximizing)} returning the minimax value of a
game tree \emph{for the player to move}. A game tree is a nested structure: an
internal node is a non-empty list of child subtrees; a leaf is a number (the
payoff from the maximizing player's viewpoint). If \texttt{maximizing} is
\texttt{True}, return the max over children (each evaluated with
\texttt{maximizing=False}); otherwise return the min (children with
\texttt{maximizing=True}).

\textbf{Example.} For the tree \texttt{[[3,5],[2,9]]} with the root maximizing:
the two MIN children evaluate to $3$ and $2$, and the MAX root returns $3$.

\subsection*{Problem 3 -- Alpha-Beta Pruning}

Implement \texttt{alphabeta(node, maximizing, alpha, beta)} returning
\texttt{(value, leaves\_visited)}. Carry the bounds $\alpha$ (best value the
maximizer can guarantee) and $\beta$ (best the minimizer can guarantee) down the
recursion, updating them after each child and \emph{breaking} out of the loop
when $\alpha \ge \beta$.

\textbf{Verify} that \texttt{alphabeta} returns the \emph{same value} as plain
minimax (\texttt{minimax\_count}, provided) on every test tree, while visiting
\emph{no more} leaves -- and strictly fewer leaves in total across all the test
trees. This demonstrates that pruning speeds up the search without changing the
answer.

\newpage
\section{Fixed-Parameter Tractability and Special Structure}

\subsection*{Problem 4 -- Vertex Cover via Bounded-Search-Tree FPT}

Implement \texttt{vertex\_cover\_fpt(n, edges, k)} deciding whether the graph
has a vertex cover of size $\le k$, returning the cover (a set) or \texttt{None}.
Use the bounded search tree: while uncovered edges remain and the budget is
positive, pick any uncovered edge $(u,v)$ and branch on adding $u$ (recurse with
budget $k-1$) and on adding $v$. The tree has depth $\le k$ and branching factor
$2$, giving $O(2^k \cdot (n+m))$ time -- FPT in $k$.

\textbf{Example.} On a triangle (edges $(0,1),(1,2),(0,2)$),
\texttt{vertex\_cover\_fpt(3, ..., 1)} returns \texttt{None} but
\texttt{vertex\_cover\_fpt(3, ..., 2)} returns a valid $2$-vertex cover.

\subsection*{Problem 5 -- Verify FPT Vertex Cover vs Brute Force}

Implement \texttt{verify\_fpt\_vertex\_cover(num\_trials, n, seed)}: for each of
\texttt{num\_trials} random graphs (\texttt{generate\_graph(n, p=0.4, seed=seed
+ trial)}), find the minimum vertex cover size via both
\texttt{brute\_force\_vertex\_cover} (\texttt{starter\_code.py}) and the FPT
decider from Problem 4 (called with increasing $k$). Return \texttt{True} iff
the sizes agree \emph{and} the FPT threshold is exact (a cover of size
\texttt{opt} exists but none of size \texttt{opt}$-1$), for every trial. Also
implement the helper \texttt{min\_vertex\_cover\_via\_fpt(n, edges)}.

\subsection*{Problem 6 -- Kernelization for Vertex Cover}

Implement \texttt{kernelize\_vertex\_cover(n, edges, k)} applying two reduction
rules exhaustively for the parameterized question ``is there a cover of size
$\le k$?'':
\begin{verbatim}
Rule 1 (high degree): any vertex of degree > k MUST be in every cover
    of size <= k. Force it in, delete it and its edges, decrement k.
Rule 2 (isolated): a degree-0 vertex covers nothing -- drop it.
\end{verbatim}
Return \texttt{(remaining\_vertices, remaining\_edges, k\_reduced, forced)},
where \texttt{forced} is the set of vertices forced into the cover by Rule 1.
The original instance is a yes-instance iff the reduced one is (with budget
\texttt{k\_reduced}).

\textbf{Example.} For a star with center $0$ and $5$ leaves with $k=1$: vertex
$0$ has degree $5 > 1$, so it is forced; its edges vanish and the leaves become
isolated and are dropped, leaving an empty instance with $k=0$.

\subsection*{Problem 7 -- Maximum Independent Set on a Tree via DP}

Implement \texttt{max\_independent\_set\_tree(n, edges)} returning a maximum
independent set of a tree in linear time, via a post-order DP rooted at vertex
$0$:
\begin{verbatim}
incl[v] = 1 + sum over children c of excl[c]
excl[v] = sum over children c of max(incl[c], excl[c])
\end{verbatim}
The optimum size is \texttt{max(incl[root], excl[root])}; reconstruct the actual
set by walking down from the root (if a vertex is included, all its children are
excluded). Also implement \texttt{verify\_tree\_mis(num\_trials, n, seed)}
confirming the DP size matches \texttt{brute\_force\_independent\_set} on random
trees (\texttt{generate\_tree}).

\textbf{Example.} On the path $0$--$1$--$2$--$3$--$4$ the maximum independent set
has size $3$ (e.g.\ $\{0,2,4\}$).

\subsection*{Problem 10 -- 2-SAT in Polynomial Time (Implication Graph + SCC)}

Although 3-SAT is NP-complete, 2-SAT is solvable in linear time. Implement
\texttt{two\_sat(num\_vars, clauses)} returning a satisfying assignment
\texttt{\{var: bool\}} or \texttt{None} if the 2-CNF formula is unsatisfiable.
Build the implication graph (each clause $(a \lor b)$ gives $\lnot a \Rightarrow
b$ and $\lnot b \Rightarrow a$), compute strongly connected components (Kosaraju
or Tarjan), and report unsatisfiable iff some variable $x$ has $x$ and $\lnot x$
in the same SCC. Otherwise read off an assignment from the SCC order.

\textbf{Verify} both directions against a brute-force satisfiability check over
all $2^n$ assignments on random formulas (\texttt{generate\_2sat}); use
\texttt{cnf\_is\_satisfied} (\texttt{starter\_code.py}) to check assignments.

\subsection*{Problem 12 -- Set Cover via Bitmask / FPT Search}

Implement \texttt{set\_cover\_fpt(universe, subsets, k)} returning the indices of
at most $k$ subsets whose union is \texttt{universe} (the smallest found), or
\texttt{None} if no cover of size $\le k$ exists. Use a bounded search tree
parameterized by $k$: while elements remain uncovered and budget $> 0$, pick any
uncovered element and branch over the subsets containing it. Represent covered
elements as a \emph{bitmask} for fast union and coverage tests. Also implement
\texttt{min\_set\_cover\_via\_fpt(universe, subsets)} (try $k=0,1,\dots$) and
verify against \texttt{brute\_force\_set\_cover} (provided) on small random
instances.

\newpage
\section{Branch and Bound and State-Space Search}

\subsection*{Problem 8 -- Branch-and-Bound 0/1 Knapsack}

Implement \texttt{knapsack\_branch\_and\_bound(values, weights, capacity)}
returning the maximum total value of a 0/1 knapsack via branch and bound. Sort
items by value/weight ratio (descending) for a tight relaxation, explore the
include/exclude binary tree, and at each node compute the \emph{fractional}
(LP-relaxation) upper bound (\texttt{fractional\_bound}, provided); prune a
subtree when its bound does not exceed the best value found so far.

\textbf{Verify} that the optimum equals the standard $O(nW)$ DP
(\texttt{dp\_knapsack}, provided) on a fixed instance and on $200$ random
instances.

\subsection*{Problem 9 -- Branch-and-Bound TSP}

Implement \texttt{tsp\_branch\_and\_bound(dist)} returning the length of the
shortest Hamiltonian cycle through all cities, starting and ending at city $0$,
for a symmetric distance matrix \texttt{dist}. Branch over the order in which
cities are visited and prune a partial tour when \texttt{cost\_so\_far +
lower\_bound} $\ge$ the best complete tour found so far. The provided
\texttt{\_lower\_bound} adds, for the current city and each unvisited city, the
cheapest edge that could still leave it -- an admissible (never-overestimating)
bound.

\textbf{Verify} that the optimum matches Held--Karp DP (\texttt{held\_karp},
provided) on random instances with small $n$.

\subsection*{Problem 11 -- Iterative-Deepening DFS on a State Space}

Implement \texttt{iddfs(start, goal, side, max\_depth)} returning a
\emph{shortest} solution path (a list of states from \texttt{start} to
\texttt{goal} inclusive) for a sliding puzzle, or \texttt{None} if no solution
within \texttt{max\_depth} moves. A state is a flattened tuple of a
\texttt{side}$\times$\texttt{side} board with $0$ as the blank;
\texttt{neighbors(state, side)} (provided) returns the states reachable by one
slide. Run depth-limited DFS with limit $1, 2, 3, \dots$ until the goal is
found; the first successful limit yields a shortest path. Track visited states
on the current path to avoid cycles.

\textbf{Example.} On the $2\times2$ puzzle, one slide from the goal returns a
path of length $2$ (start, goal).

\newpage
\section{LeetCode Practice}

After completing the problems above, practise this week's techniques (game
search / minimax, FPT / bitmask search, state-space search, branch-and-bound) on
the following \href{https://leetcode.com}{LeetCode} problems. Difficulty is
LeetCode's own label.

\begin{itemize}
  \item \href{https://leetcode.com/problems/word-ladder/}{Word Ladder} -- \emph{Hard} -- State-space BFS.
  \item \href{https://leetcode.com/problems/sliding-puzzle/}{Sliding Puzzle} -- \emph{Hard} -- State-space search.
  \item \href{https://leetcode.com/problems/smallest-sufficient-team/}{Smallest Sufficient Team} -- \emph{Hard} -- Set cover (FPT / bitmask).
  \item \href{https://leetcode.com/problems/stone-game/}{Stone Game} -- \emph{Medium} -- Game / minimax DP.
  \item \href{https://leetcode.com/problems/stone-game-ii/}{Stone Game II} -- \emph{Medium} -- Game / minimax DP.
  \item \href{https://leetcode.com/problems/can-i-win/}{Can I Win} -- \emph{Medium} -- Game search + memo.
  \item \href{https://leetcode.com/problems/cat-and-mouse/}{Cat and Mouse} -- \emph{Hard} -- PSPACE game (BFS on states).
  \item \href{https://leetcode.com/problems/maximum-students-taking-exam/}{Maximum Students Taking Exam} -- \emph{Hard} -- Bitmask FPT.
\end{itemize}

\end{document}
